%% This is a JAIR Example File Compiled by Nicholas Mattei (nsmattei@tulane.edu) 
%% and Odd Erik Gundersen (odderik@ntnu.no)
%% and Mykel Kochenderfer (mykel@stanford.edu)
%% September 15, 2025
%%
%% This file is based off the ACM Latex Template https://www.acm.org/publications/proceedings-template
%% Revision 2.12 (12/28/2024)
%% 
%% Please see https://www.jair.org/index.php/jair for more information and submission instructions.
%%

%% The first command in your LaTeX source must be the \documentclass
%% command.
%%
%% For final camera ready please change the
%% command to \documentclass[]{jair}.
%%
\documentclass[manuscript, screen, review]{jair}


\setcopyright{cc}
\copyrightyear{2025}
\acmDOI{10.1613/jair.1.xxxxx}

%%
\JAIRAE{Insert JAIR AE Name}
\JAIRTrack{} % Insert JAIR Track Name only if part of a special track
\acmVolume{1}
\acmArticle{1}
\acmMonth{1}
\acmYear{2026}

\RequirePackage[
  datamodel=acmdatamodel,
  style=acmauthoryear,
  backend=biber,
  giveninits=true,
  uniquename=init
  ]{biblatex}

\renewcommand*{\bibopenbracket}{(}
\renewcommand*{\bibclosebracket}{)}

%%
%% For managing citations, use BibLaTeX with the acmauthoryear style.
%% The next line specifies the bibliography file.
\addbibresource{bibliography.bib}

%% Additional packages for paper generator
\usepackage[utf8]{inputenx}
\graphicspath{{images/}}
\DeclareGraphicsExtensions{.pdf,.png,.jpeg,.jpg}
%% Note: hyperref and xcolor are already loaded by acmart.cls

%%
%% end of the preamble, start of the body of the document source.
\begin{document}

%%
%% The "title" command has an optional parameter,
%% allowing the author to define a "short title" to be used in page headers.
%% Usage: \title[short title]{full title}
%%
\title{%%TITLE%%}

%%
%% The "author" command and its associated commands are used to define
%% the authors and their affiliations.
%% Of note is the shared affiliation of the first two authors, and the
%% "authornote" and "authornotemark" commands
%% used to denote shared contribution to the research and/or corresponding author.
%%
%% Each author must be defined separately for accurate metadata identification.
%% As an exception, multiple authors may share one affiliation. 
%% Authors' names should not be abbreviated; use full first names wherever possible. 
%% Include authors' e-mail addresses whenever possible. 
%% Including ORCID information for each author is highly recommended.
%% Example ORCID: \orcid{0000-0002-2037-3694}
%%
%% Grouping authors' names or e-mail addresses is not acceptable:
%%   BAD: \author{Brooke Aster, David Mehldau}
%%   BAD: \email{dave,judy,steve@university.edu}
%%
%% The authornote and authornotemark commands allow a note to apply to 
%% multiple authors -- for example, if the first two authors contributed equally.
%%
%% Author format:
%% \author{Author Name}
%% \authornote{Corresponding Author.}  % optional
%% \orcid{0000-0000-0000-0000}  % optional
%% \email{author@institution.edu}
%% \affiliation{%
%%   \institution{Institution Name}
%%   \city{City}
%%   \state{State/Province}
%%   \country{Country}
%% }
%%

%%BEGIN_AUTHOR%%
\author{{{name}}}
\email{{{email}}}
\affiliation{%
  \institution{{{affiliation}}}
  \city{{{city}}}
  \country{{{country}}}
}
%%END_AUTHOR%%

%% The short list of authors must be made of the list of all authors' lastnames.
%% If this is too long and overlaps other information printed in the page headers, use:
%% \renewcommand{\shortauthors}{First Author et al.}
\renewcommand{\shortauthors}{%%SHORTAUTHORS%%}

%%
%% The abstract is a short summary of the work to be presented in the
%% article.
\begin{abstract}
%%ABSTRACT%%
\end{abstract}

%% JAIR Note: 
%% Do not include ACM CCS Concepts or Keywords

%% To be updated by authors.
\received{Date received}
\received[accepted]{Date accepted}

%% This command processes the author and affiliation and title
%% information and builds the first part of the formatted document.
\maketitle

\section{Introduction}
%%INTRODUCTION%%

\section{Related Work}
%%RELATED_WORK%%

\section{Methods}
%%METHODS%%

\section{Results}
%%RESULTS%%

\section{Discussion}
%%DISCUSSION%%

\section{Conclusion}
%%CONCLUSION%%

%%
%% The acknowledgments section is defined using the "acks" environment
%% (and NOT an unnumbered section). This ensures the proper
%% identification of the section in the article metadata, and the
%% consistent spelling of the heading.
\begin{acks}
%%ACKNOWLEDGEMENTS%%
\end{acks}

%%
%% The next line prints the references.
\printbibliography


%%
%% If your work has an appendix, this is the place to put it.
%% Start the appendix with the "appendix" command:
%%   \appendix
%% Note that in the appendix, sections are lettered, not numbered.


\end{document}
\endinput
%%
%% End of file `sample-acmlarge.tex'.

%% ===========================================================================
%% JAIR TEMPLATE DOCUMENTATION (preserved as comments for reference)
%% ===========================================================================
%%
%% MODIFICATIONS:
%% Modifying the template --- including but not limited to: adjusting
%% margins, typeface sizes, line spacing, paragraph and list definitions,
%% and the use of the \vspace command to manually adjust the
%% vertical spacing between elements of your work --- is not allowed.
%% Your document will be returned to you for revision if modifications 
%% are discovered.
%%
%% ---------------------------------------------------------------------------
%% TYPEFACES:
%% The "acmart" document class requires the use of the
%% "Libertine" typeface family. Your TeX installation should include
%% this set of packages. Please do not substitute other typefaces. The
%% "lmodern" and "ltimes" packages should not be used,
%% as they will override the built-in typeface families.
%%
%% ---------------------------------------------------------------------------
%% TITLE INFORMATION:
%% The title of your work should use capital letters appropriately -
%% https://capitalizemytitle.com/ has useful rules for capitalization. 
%% Use the title command to define the title of your work. 
%% If your work has a subtitle, define it with the subtitle command. 
%% Do not insert line breaks in your title.
%%
%% If your title is lengthy, you must define a short version to be used
%% in the page headers, to prevent overlapping text:
%%   \title[short title]{full title}
%%
%% ---------------------------------------------------------------------------
%% RIGHTS INFORMATION:
%% JAIR articles are freely available, and a copyright notice referencing 
%% the Creative Commons CC BY license is included in the template. 
%% This copyright notice should be included in all JAIR articles.
%%
%% ---------------------------------------------------------------------------
%% SECTIONING COMMANDS:
%% Your work should use standard LaTeX sectioning commands:
%% section, subsection, subsubsection, and paragraph. They should be numbered; 
%% do not remove the numbering from the commands.
%%
%% Simulating a sectioning command by setting the first word or words of
%% a paragraph in boldface or italicized text is not allowed.
%%
%% ---------------------------------------------------------------------------
%% TABLES:
%% The "acmart" document class includes the "booktabs"
%% package --- https://ctan.org/pkg/booktabs --- for preparing high-quality tables.
%%
%% Table captions are placed above the table.
%%
%% Because tables cannot be split across pages, the best placement for
%% them is typically the top of the page nearest their initial cite. To
%% ensure this proper "floating" placement of tables, use the
%% environment table to enclose the table's contents and the
%% table caption. The contents of the table itself must go in the
%% tabular environment.
%%
%% Always use midrule to separate table header rows from data rows, and
%% use it only for this purpose. This enables assistive technologies to
%% recognise table headers and support their users in navigating tables.
%%
%% Example table:
%% \begin{table}
%%   \caption{Frequency of Special Characters}
%%   \label{tab:freq}
%%   \begin{tabular}{@{}ccl@{}}
%%     \toprule
%%     Non-English or Math&Frequency&Comments\\
%%     \midrule
%%     \O & 1 in 1,000& For Swedish names\\
%%     $\pi$ & 1 in 5& Common in math\\
%%     \$ & 4 in 5 & Used in business\\
%%     $\Psi^2_1$ & 1 in 40,000& Unexplained usage\\
%%   \bottomrule
%% \end{tabular}
%% \end{table}
%%
%% ---------------------------------------------------------------------------
%% MATH EQUATIONS:
%% You may want to display math equations in three distinct styles:
%% inline, numbered or non-numbered display. 
%%
%% JAIR authors are discouraged from using sophisticated mathematical notation 
%% in the abstract (especially notation that is difficult to reproduce in 
%% vanilla HTML), since the abstract will also be formatted in HTML on the 
%% JAIR web site.
%%
%% Inline (In-text) Equations:
%% A formula that appears in the running text is called an inline or
%% in-text formula. It is produced by the math environment,
%% with the usual \begin...\end construction or with \$...\$.
%%
%% Display Equations:
%% A numbered display equation is produced by the equation environment.
%% An unnumbered display equation is produced by the displaymath environment.
%%
%% ---------------------------------------------------------------------------
%% FIGURES:
%% The "figure" environment should be used for figures. One or
%% more images can be placed within a figure. If your figure contains
%% third-party material, you must clearly identify it as such.
%%
%% Your figures should contain a caption which describes the figure to the reader.
%% Figure captions are placed below the figure.
%%
%% Every figure should also have a figure description. These descriptions 
%% convey what is in the image to someone who cannot see it. They are also 
%% used by search engine crawlers for indexing images, and when images 
%% cannot be loaded.
%%
%% A figure description must be unformatted plain text less than 2000
%% characters long (including spaces). Figure descriptions
%% should not repeat the figure caption – their purpose is to capture
%% important information that is not already provided in the caption or
%% the main text of the paper.
%%
%% For additional information regarding how best to write figure descriptions 
%% and why doing this is so important, please see
%% https://www.acm.org/publications/taps/describing-figures/
%%
%% Example figure:
%% \begin{figure}[ht]
%%   \centering
%%   \includegraphics[width=0.6\linewidth]{image-file}
%%   \caption{Caption text here.}
%%   \Description{Description for accessibility.}
%% \end{figure}
%%
%% ---------------------------------------------------------------------------
%% CITATIONS AND BIBLIOGRAPHIES:
%% For managing citations, use bibliography files in BibTeX format.
%% You can then use BibLaTeX with the acmauthoryear style. 
%% Do not use the acmnumeric style.
%%
%% The bibliography is included in your source document with the command:
%%   \addbibresource{bibfile.bib}
%% where "bibfile.bib" is the name of the bibliography file.
%%
%% We can reference the work of the authors using the \citet command. 
%% Do not use a parenthetical citation as a noun in a sentence.
%%
%% ---------------------------------------------------------------------------
%% ACKNOWLEDGMENTS:
%% Identification of funding sources and other support, and thanks to
%% individuals and groups that assisted in the research and the
%% preparation of the work should be included in an acknowledgment
%% section, which is placed just before the reference section.
%%
%% This section has a special environment:
%%   \begin{acks}
%%   ...
%%   \end{acks}
%% so that the information contained therein can be more easily collected
%% during the article metadata extraction phase, and to ensure
%% consistency in the spelling of the section heading.
%%
%% Authors should not prepare this section as a numbered or unnumbered section;
%% please use the "acks" environment.
%%
%% ---------------------------------------------------------------------------
%% APPENDICES:
%% If your work needs an appendix, add it before the
%% \end{document} command at the conclusion of your source document.
%%
%% Start the appendix with the "appendix" command:
%%   \appendix
%% and note that in the appendix, sections are lettered, not numbered.
%%
%% ---------------------------------------------------------------------------
%% REPRODUCIBILITY CHECKLIST FOR JAIR:
%% Select the answers that apply to your research -- one per item.
%%
%% ALL ARTICLES:
%% 1. All claims investigated in this work are clearly stated. [yes/partially/no]
%% 2. Clear explanations are given how the work reported substantiates the claims. [yes/partially/no]
%% 3. Limitations or technical assumptions are stated clearly and explicitly. [yes/partially/no]
%% 4. Conceptual outlines and/or pseudo-code descriptions of the AI methods introduced in this work are provided, and important implementation details are discussed. [yes/partially/no/NA]
%% 5. Motivation is provided for all design choices, including algorithms, implementation choices, parameters, data sets and experimental protocols beyond metrics. [yes/partially/no]
%%
%% ARTICLES CONTAINING THEORETICAL CONTRIBUTIONS:
%% Does this paper make theoretical contributions? [yes/no]
%% If yes:
%% 1. All assumptions and restrictions are stated clearly and formally. [yes/partially/no]
%% 2. All novel claims are stated formally (e.g., in theorem statements). [yes/partially/no]
%% 3. Proofs of all non-trivial claims are provided in sufficient detail. [yes/partially/no]
%% 4. Complex formalism, such as definitions or proofs, is motivated and explained clearly. [yes/partially/no]
%% 5. The use of mathematical notation and formalism serves the purpose of enhancing clarity and precision. [yes/partially/no]
%% 6. Appropriate citations are given for all non-trivial theoretical tools and techniques. [yes/partially/no]
%%
%% ARTICLES REPORTING ON COMPUTATIONAL EXPERIMENTS:
%% Does this paper include computational experiments? [yes/no]
%% If yes:
%% 1. All source code required for conducting experiments is included online. [yes/partially/no]
%% 2. The source code comes with a license for reproducibility purposes. [yes/partially/no]
%% 3. The source code comes with a license for research purposes. [yes/partially/no]
%% 4. Raw, unaggregated data from all experiments is included online. [yes/partially/no]
%% 5. The unaggregated data comes with a license for reproducibility purposes. [yes/partially/no]
%% 6. The unaggregated data comes with a license for research purposes. [yes/partially/no]
%% 7. If an algorithm depends on randomness, then the method for random numbers is described. [yes/partially/no/NA]
%% 8. The execution environment for experiments is described. [yes/partially/no]
%% 9. The evaluation metrics used in experiments are clearly explained. [yes/partially/no]
%% 10. The number of algorithm runs used to compute each result is reported. [yes/no]
%% 11. Reported results have not been "cherry-picked". [yes/partially/no]
%% 12. Analysis of results goes beyond single-dimensional summaries. [yes/no]
%% 13. All (hyper-) parameter settings have been reported. [yes/partially/no/NA]
%% 14. The number and range of (hyper-) parameter settings explored is indicated. [yes/partially/no/NA]
%% 15. Appropriately chosen statistical hypothesis tests are used. [yes/partially/no/NA]
%%
%% ARTICLES USING DATA SETS:
%% Does this work rely on one or more data sets? [yes/no]
%% If yes:
%% 1. All newly introduced data sets are included online. [yes/partially/no/NA]
%% 2. The newly introduced data set comes with a license for reproducibility. [yes/partially/no]
%% 3. The newly introduced data set comes with a license for research. [yes/partially/no]
%% 4. All data sets drawn from the literature are accompanied by appropriate citations. [yes/no/NA]
%% 5. All data sets drawn from the existing literature are publicly available. [yes/partially/no/NA]
%% 6. All new data sets and non-public data sets are described in detail. [yes/partially/no/NA]
%% 7. All methods used for preprocessing, augmenting, batching or splitting data sets are described. [yes/partially/no/NA]
%%
